\begin{proof}
仅对多项式情形给出证明，三角情形类似。

\textbf{充分性}：设 $x_1<\cdots<x_{n+2}$ 是 $p$ 的一组交错偏离点，即
\[
|f(x_j)-p(x_j)| = e_n(f),\quad j=1,\dots,n+2,
\]
且相邻点误差符号交替。不妨设 $f(x_j)-p(x_j) = (-1)^{j-1}e_n(f)$。
若 $p$ 不是最佳逼近，则存在 $q\in P_n$ 使 $\|f-q\| < \|f-p\|$。
考虑 $d(x)=p(x)-q(x)\in P_n$。在 $x_j$ 处，
\[
d(x_j) = (p(x_j)-f(x_j)) - (q(x_j)-f(x_j)).
\]
因 $\|f-q\|<\|f-p\|$，故 $d(x_j)$ 的符号与 $p(x_j)-f(x_j)$ 相同。于是 $d$ 在 $x_j$ 处依次变号，从而 $d$ 至少有 $n+1$ 个零点，但 $d$ 是 $n$ 次多项式，于是 $d$ 至少有$ n+1$ 个零点。
但 $d$ 为次数不超过 $n$ 的多项式，
除非 $d\equiv0$，
否则不可能有 $n+1$ 个不同零点。
故 $d\equiv0$，即 $p\equiv q$，
与假设矛盾。

\textbf{必要性}：
\begin{enumerate}
    \item \textbf{交错偏离点有限且小于 $n+2$} 

若交错偏离点个数 $N < n+2$，但 $p$ 仍是最佳逼近。记 $r(x)=f(x)-p(x)$，$M=\|r\|$。设
\[
a<x_1<\cdots<x_N<b
\]
为全体交错偏离点。

对每个 $i=1,\dots,N-1$，
由连续性可知存在
\[
\xi_i\in(x_i,x_{i+1})
\]
使得
\[
r(\xi_i)=0.
\]

再记
\[
\xi_0=a,\qquad \xi_N=b.
\]将 $[a,b]$ 用 $$\xi_1,\dots,\xi_{N-1}$$ 分割成若干子区间，则 $r$ 在每个子区间上满足
\[
-M+\alpha < r(x)  \leqslant M \quad\text{或}\quad -M  \leqslant r(x) < M-\alpha.
\]
因为我们可以假设在区间 $[a, \xi_1]$ 上有 $r(x)  \leqslant 0$。因为 $\|r\| = M$，从而 $|r(x)|  \leqslant  M$，从而此时 $r(x) \geqslant -M$ 是成立的。

由闭区间上连续函数的性质，存在 $0 < \alpha_1 < M $ 使得 $-M < r(x) < M -\alpha_1$ 在 $[a, \xi_1]$ 上成立。同理在每个子区间上都存在 $ 0 < \alpha_i < M $ 使得上述不等式成立。取 $\alpha = \min\{\alpha_i\}$，则在每个子区间上都有
\[-M   \leqslant  r(x) < M - \alpha.
\]或者
\[-M + \alpha < r(x)  \leqslant  M \]成立。

由于 $r$ 在每个子区间上不变号，
且 $\phi$ 在穿过每个 $\xi_i$ 后也恰变号一次，
故 $r$ 与 $\phi$ 在每个子区间上同号。

令 $\phi(x) = \prod_{i=1}^{N-1}(x-\xi_i)\in P_{N-1}\subseteq P_n$。构造 $q(x)=p(x)+c\phi(x)$。我们通过选取合适的 $c$ 去导出 $q(x)$ 是比 $p(x)$ 更好的逼近，从而导致矛盾。

不妨设在 $[a, \xi_1]$ 上 $r(x) \leqslant  0$ 并且 $\phi(x)$ 在 $[a, \xi_1]$ 上满足 $\phi(x)  \leqslant  0$（这只是符号上的问题）。
记 $K = \|\phi\|$.

此时 $-M  \leqslant  r(x) < M-\alpha $  注意到 $\phi(x)$ 与 $r(x)$ 在 $[a, \xi_1]$ 上符号相同。因为
$$
\begin{aligned}
f(x) - q(x) &= f(x) - p(x) - c\phi(x)  \\
              &= r(x) - c\phi(x) \\
              &= r(x) + c|\phi(x)| \\
\end{aligned}$$
这里由于 $r(x)  \leqslant  0$ 并且 $\phi(x)  \leqslant  0$，所以 
当 $c$ 满足 $ 0 < c| \phi(x) |  \leqslant  |c\cdot K| < \frac{\alpha}{2}$ 并且 $c>0$ 时
有
$$
\begin{aligned}
-M  < r(x) - c\phi(x) < M -\alpha + \frac{\alpha}{2} \\ 
\end{aligned}
$$
同理，在区间 $[\xi_1, \xi_2]$ 上我们能够利用
$$
\begin{aligned}
 -M +\alpha < r(x)  \leqslant  M
\end{aligned}
$$
和
$$
\begin{aligned}
f(x) - q(x) &= f(x) - p(x) - c\phi(x)  \\
              &= r(x) - c\phi(x) \\
              &= r(x) - c|\phi(x)| \\
\end{aligned}
$$
得到
$$
\begin{aligned}
-M + \alpha - \frac{\alpha}{2} < r(x) - c\phi(x) < M 
\end{aligned}$$
注意到上述过程与子区间的选择无关，因此在每个子区间上都满足上述不等式。
从而有 $$\|f-q\| < M = \|f-p\|$$，并且根据 $q$ 的构造可以得知它是不超过$n$ 次的多项式，
于是与 $\|f-q\| < \|f-p\|$矛盾。

\item \textbf{交错偏离点有限且大于等于 $n+2$} 

由于交错偏离点按符号交替出现，
故可从中依次选取
$n+2$
个点构成一组交错点。
由充分性部分，
$p$
即为最佳逼近。

\item \textbf{交错偏离点有无穷个}记若交错偏离点有无穷多个。记
$$
E_+=\{x\in[a,b]:r(x)=M\},\qquad
E_-=\{x\in[a,b]:r(x)=-M\}.
$$

由于 $r$ 连续，而 $\{M\},\{-M\}$ 为闭集，因此
$$
E_+=r^{-1}(\{M\}),\qquad
E_-=r^{-1}(\{-M\})
$$
均为闭集，从而为紧集。

由于交错偏离点无穷多个，因此在区间 $[a,b]$ 中，$E_+$ 与 $E_-$ 必无限次交替出现。

先取
$$
x_1=\min(E_+\cup E_-).
$$
不妨设
$$
r(x_1)=M.
$$

由于交错偏离点无穷多个，因此集合
$$
\{x>x_1:x\in E_-\}
$$
非空。又因为 $E_-$ 为闭集，所以该集合在 $[a,b]$ 中取到最小元。定义
$$
x_2=\min\{x>x_1:x\in E_-\}.
$$
于是
$$
x_1<x_2,\qquad r(x_2)=-M.
$$

同理，由于交错偏离点无穷多个，集合
$$
\{x>x_2:x\in E_+\}
$$
非空，并且取到最小元。定义
$$
x_3=\min\{x>x_2:x\in E_+\}.
$$
于是
$$
x_2<x_3,\qquad r(x_3)=M.
$$

继续递归构造：
若已定义
$$
x_k,
$$
则定义
$$
x_{k+1}
=
\min\{x>x_k:x\in E_{\sigma_{k+1}}\},
$$
其中
$$
\sigma_{k+1}=
\begin{cases}
+,& r(x_k)=-M,\\
-,& r(x_k)=M.
\end{cases}
$$

由于交错偏离点无穷多个，上述过程不会在有限步停止。

因此可得到
$$
x_1<x_2<\cdots<x_{n+2}
$$
满足
$$
r(x_i)=(-1)^iM.
$$

故已经存在不少于 $n+2$ 个交错偏离点。
\end{enumerate}
\end{proof}
